\documentclass[11pt]{article}

\usepackage[margin=1.1in]{geometry}
\usepackage{amsmath,amssymb,amsthm,mathtools}
\usepackage{enumitem}
\usepackage{hyperref}
\usepackage[nameinlink,noabbrev]{cleveref}

\hypersetup{colorlinks=true,linkcolor=blue,citecolor=blue,urlcolor=blue}

\newtheorem{theorem}{Theorem}
\newtheorem{lemma}[theorem]{Lemma}
\newtheorem{proposition}[theorem]{Proposition}
\newtheorem{corollary}[theorem]{Corollary}
\theoremstyle{remark}
\newtheorem{remark}[theorem]{Remark}
\theoremstyle{plain}
\newtheorem{thmz}{Theorem}
\renewcommand{\thethmz}{Z}
\crefname{thmz}{Theorem}{Theorems}

\DeclareMathOperator{\supp}{supp}
\providecommand{\doi}[1]{\href{https://doi.org/#1}{doi:#1}}
\newcommand{\tup}{t_{\mathrm{up}}}

\title{Stationary robustness of random walks against a strong adversary at large degree}
\author{Yitzchak Shmalo\\
Einstein Institute of Mathematics, The Hebrew University of Jerusalem\\
Givat Ram, Jerusalem, Israel}
\date{}

\begin{document}
\maketitle

\begin{abstract}
A strong $\varepsilon$-adversary for a Markov chain replaces each transition, with probability $\varepsilon$, by a step
of its own choosing inside the support of the chain. The upper stationary probability $\pi^+(S)$ of a set $S$ is the
largest long-run fraction of time in $S$ that such an adversary can force. For the simple random walk on a graph
with degrees at most $d$, Zuckerman stated $\pi^+(S)=O(\pi(S)^{1-o(1)})$ for $\varepsilon$ of order
$1/\sqrt{d\,\tup}$, where $\tup$ is a variant of the mixing time, and conjectured the same with the mixing time
itself. Both fail on the complete graph $K_n$ with $\varepsilon=c/\sqrt{n-1}$, where an adversary keeps the walk at
one vertex at least a fraction $\varepsilon/2$ of the time. For a reversible chain whose positive entries are at
least $1/D$ and whose mixing time is $\tau$, we prove that $\varepsilon^2D\tau\le1/20$ implies
$\pi^+(S)\le2\pi(S)^{1/2}$ for every $S$. The exponent $1/2$ cannot be raised, even when $\varepsilon\sqrt{D\tau}$
is required to be smaller than any fixed constant.
\end{abstract}

\noindent\textbf{Keywords:} random walk, adversarial perturbation, stationary distribution, mixing time.

\noindent\textbf{MSC 2020:} 60J10, 05C81.

\section{Introduction}\label{sec:intro}

Zuckerman~\cite{Zu} asks how far an adversary that controls a small fraction of the steps of a random walk can raise
the long-run fraction of time the walk spends in a set. For walks on graphs of small degree he states a bound close to
the stationary probability of the set. At large degree that bound fails, and we prove a bound with exponent $1/2$
that holds for every reversible chain and cannot be improved in the exponent.

Let $P$ be the transition matrix of an irreducible Markov chain on a finite set $V$ with stationary distribution
$\pi$. Following Zuckerman~\cite{Zu}, a \emph{strong $\varepsilon$-adversary} chooses at each time $i$ a transition
matrix $B_i$, which may depend on the whole history $X_0,\dots,X_i$, with $\supp B_i(x,\cdot)\subseteq\supp P(x,\cdot)$
for every $x$, and the walk moves according to
\[
\Pr[X_{i+1}=y\mid X_0,\dots,X_i]=(1-\varepsilon)P(X_i,y)+\varepsilon B_i(X_i,y).
\]
Call a matrix $Q$ \emph{admissible} if $Q=(1-\varepsilon)P+\varepsilon B$ for a transition matrix $B$ with
$\supp B(x,\cdot)\subseteq\supp P(x,\cdot)$ for every $x$; at time $i$ the walk above moves by the admissible matrix
$Q_i=(1-\varepsilon)P+\varepsilon B_i$. Azar, Broder, Karlin, Linial and Phillips~\cite{ABKLP} study a controller
that fixes one such $B$ in order to help the walk, and Haslegrave, Sauerwald and Sylvester~\cite{HSS} study helpful
controllers that may vary in time; in~\cite{Zu} the controller is an adversary. Writing $N_t(S)$ for the number of
visits to $S\subseteq V$ at times $0,\dots,t-1$, Zuckerman defines the upper stationary probability
\[
\pi^+(S)=\sup\,\limsup_{t\to\infty}\frac{N_t(S)}{t},
\]
the supremum being over strong $\varepsilon$-adversaries \cite[Definition~1.1]{Zu}. How far $\pi^+(S)$ can exceed
$\pi(S)$ is the stationary robustness of the chain, in the terms of~\cite[Section~1.2]{Zu}. Let
$\tau=\min\{t:\max_x\|P^t(x,\cdot)-\pi\|_{\mathrm{TV}}\le 1/4\}$ be the mixing time. For the simple random walk on a
graph let $p_t(u,\cdot)$ be the law at time $t$ of the walk without the adversary started at $u$, let
$s_{\mathrm{up}}(t)=\max_{u,v}\bigl(p_t(u,v)/\pi(v)-1\bigr)$ and, for a constant $\alpha\in(0,1)$, let
$\tup(\alpha)=\min\{t:s_{\mathrm{up}}(t)\le\alpha\}$, the upper separation time of~\cite[Definitions~2.6 and~2.7]{Zu}.

\begin{thmz}[{Zuckerman~\cite[Theorem~1.2]{Zu}}]\label{thm:zu}
Let $G=(V,E)$ be a graph with all degrees at most $d$. For every strong $\varepsilon$-adversary for the simple random
walk on $G$, every starting distribution, and every $S\subseteq V$ with $\pi(S)\le 1-\alpha$,
\[
\pi^+(S)<(1+\alpha)\exp\Bigl(2\varepsilon\sqrt{2d\,\tup\ln(1/\pi(S))}\Bigr)\pi(S),\qquad \tup=\tup(\alpha).
\]
In particular, for $\varepsilon=O(1/\sqrt{d\,\tup})$ one has $\pi^+(S)=O(\pi(S)^{1-o(1)})$.
\end{thmz}

Zuckerman conjectures the same bound with $\tau$ in place of $\tup$ \cite[Conjecture~1.18]{Zu}. Section~1.1
of~\cite{Zu} treats the maximum degree as small or constant. Both statements fail when the degree is large. On the
complete graph $K_n$ with $n\ge5$ we have $d=n-1$, $\tau=1$, and $\tup(\alpha)=1$ for every
$\alpha\in[1/(n-1),1-1/n]$ (\Cref{lem:kn}). Fix $c\in(0,1]$ and let $\varepsilon=c/\sqrt{n-1}$. The adversary that
always moves to a fixed vertex $v$ gives, almost surely,
\[
\lim_{t\to\infty}\frac{N_t(\{v\})}{t}\ \ge\ \frac{\varepsilon}{2}=\frac{c}{2\sqrt{n-1}},
\]
while the right-hand side of \Cref{thm:zu} for $S=\{v\}$ is less than $2e^{2c\sqrt{2\ln n}}/n=n^{-1+o(1)}$, which is
smaller for every $n$ larger than a constant depending on $c$ (\Cref{thm:fail}). The same example contradicts the last
sentence of \Cref{thm:zu}, since here $\varepsilon=c/\sqrt{d\,\tup}$. The proof in~\cite{Zu} uses a bound at a fixed
time $t$. Let $p$ be the probability that the walk without the adversary is in $S$ at time $t$. Under the adversary
this probability is at most $e^{2tkd\varepsilon^2}p^{1-1/k}$ for every $k>1$ with $kd\varepsilon\le1$
\cite[Theorem~3.3]{Zu}. The proof of~\cite[Corollary~3.5]{Zu} takes $k=\sqrt{\ln(1/p)/(2td\varepsilon^2)}$. Then
$kd\varepsilon=\sqrt{d\ln(1/p)/(2t)}$, which exceeds $1$ when $d\ln(1/p)>2t$ (\Cref{rem:where}).

\begin{theorem}\label{thm:fail}
Let $n\ge 5$, let $c\in(0,1]$, let $\varepsilon=c/\sqrt{n-1}$, fix a vertex $v$ of $K_n$, and put
$a=\varepsilon+(1-\varepsilon)/(n-1)$. For the simple random walk on $K_n$ there is a strong $\varepsilon$-adversary
such that, from any start, almost surely
\[
\lim_{t\to\infty}\frac{N_t(\{v\})}{t}=\frac{a}{1+a}\ge\frac{\varepsilon}{2},
\]
and $\pi^+(\{v\})=a/(1+a)$. For every $\alpha\in[1/(n-1),1-1/n]$ one has $\tup(\alpha)=1$ and
$\pi(\{v\})\le1-\alpha$, and the right-hand side of \Cref{thm:zu} for $S=\{v\}$ is less than
$2e^{2c\sqrt{2\ln n}}/n$. Hence \Cref{thm:zu}, and its form with $\tau$ in place of $\tup$, fail for $K_n$ once
$c\,n\,e^{-2c\sqrt{2\ln n}}>4\sqrt{n-1}$, which holds for every $n$ larger than a constant depending on $c$.
\end{theorem}

For $n=65537$ and $\varepsilon=2^{-11}$ the limit equals $67583/134285311$. This is more than $5$ times the
right-hand side of \Cref{thm:zu} for every $\alpha\in[1/(n-1),1-1/n]$, and more than $8$ times it for $\alpha=1/4$.

The factor in \Cref{thm:zu} cannot hold uniformly in the degree. The exponent $1/2$ can, by the following argument.
Let $P$ be reversible with respect to $\pi$ with positive entries at least $1/D$, let $Q$ be admissible, let $(Y_i)$
be the chain $P$ started at $x$, with expectation $\mathbb E_x$, and let
$L=\prod_{i<\tau}Q(Y_i,Y_{i+1})/P(Y_i,Y_{i+1})$. Then
\[
\mathbb E_xL^2\le(1+\varepsilon^2D)^\tau\le e^{\varepsilon^2D\tau}
\]
with no condition on $D\varepsilon$. Write $\|h\|^2=\sum_x\pi(x)h(x)^2$. If $\varepsilon^2D\tau\le1/20$, the display
gives $\|(Q^\tau-P^\tau)h\|\le\frac14\|h\|$ for every function $h$ on $V$ (\Cref{lem:lr}), and reversibility gives
$\|(P^\tau-\Pi)h\|\le\frac12\|h\|$, where every row of $\Pi$ is $\pi$ (\Cref{lem:gap}). Let $\nu$ be a stationary
distribution of $Q$, let $f=\nu/\pi$, and let $\mathbf 1$ be the constant function $1$. Then
\[
\|f-\mathbf 1\|\le\tfrac14\|f\|+\tfrac12\|f-\mathbf 1\|,\qquad\text{so}\qquad\|f\|\le2\quad\text{and}\quad
\nu(S)\le\|f\|\,\pi(S)^{1/2}\le2\pi(S)^{1/2}
\]
(\Cref{prop:stat}). Almost surely every limit point of the empirical measure of the walk under a strong
$\varepsilon$-adversary is stationary for an admissible matrix (\Cref{lem:limit}), and this gives the following.

\begin{theorem}\label{thm:half}
Let $P$ be an irreducible transition matrix on a finite set $V$, reversible with respect to $\pi$, such that
$P(x,y)\ge 1/D$ whenever $P(x,y)>0$, and suppose $\tau<\infty$. If $\varepsilon^2D\tau\le 1/20$, then for every
strong $\varepsilon$-adversary, every starting distribution and every $S\subseteq V$, almost surely
\[
\limsup_{t\to\infty}\frac{N_t(S)}{t}\le 2\pi(S)^{1/2}.
\]
In particular $\pi^+(S)\le 2\pi(S)^{1/2}$.
\end{theorem}

For the simple random walk on a connected non-bipartite graph with degrees at most $d$ one may take $D=d$. When
$2d\varepsilon\le1$, a bound with exponent $1/2$ at a fixed time is \cite[Theorem~3.3]{Zu} with $k=2$;
\Cref{thm:half} adds the range $d\varepsilon>1/2$, the mixing time in place of $\tup$, and reversible chains. By
\Cref{thm:fail} the exponent $1/2$ cannot be raised.

\begin{corollary}\label{cor:exp}
Let $\kappa\in(0,1]$. There are constants $c_\kappa,C_\kappa>0$ such that $\pi^+(S)\le C_\kappa\pi(S)^\kappa$ for
every chain as in \Cref{thm:half}, every $\varepsilon\le c_\kappa/\sqrt{D\tau}$ and every $S\subseteq V$, if and only
if $\kappa\le1/2$.
\end{corollary}

The ``only if'' half also holds with the lazy walk on $K_n$ in place of the simple random walk (\Cref{sec:fail}).
For graphs of bounded degree \Cref{thm:zu}, where it holds, gives $\pi(S)^{1-o(1)}$, which for small $\pi(S)$ is
much smaller than the bound $2\pi(S)^{1/2}$ of \Cref{thm:half}. Whether \Cref{thm:zu} holds for bounded $d$ is open: the choice of $k$ in the proof in~\cite{Zu} is
allowed only when $d\ln(1/p)\le2t$.

On the rook's graph $K_q\square K_q$ with $q\ge41$ and $\varepsilon=1/\sqrt{20d\tau}$, \Cref{thm:half} gives
$\pi^+(\{x\})\le2/q$ at every vertex $x$. There the bound of \Cref{prop:tv} below is larger than $2/q$, the bound
of~\cite[Proposition~4.7]{HSS} is larger than $1/4$, and the bound of~\cite[Proposition~4.8]{HSS} is larger than $1$
(\Cref{prop:rook}).

\section{Empirical measures}\label{sec:emp}

For the walk under a strong $\varepsilon$-adversary, with the matrices $Q_i=(1-\varepsilon)P+\varepsilon B_i$, put
\[
\mu_t(x)=\frac1t\sum_{i<t}\mathbf 1\{X_i=x\},\qquad F_t(x,y)=\frac1t\sum_{i<t}\mathbf 1\{X_i=x\}Q_i(x,y).
\]
The first lemma is the strong law for bounded martingale differences; we include the short proof.

\begin{lemma}\label{lem:emp}
Almost surely, $\mu_t(y)-\sum_xF_t(x,y)\to0$ for every $y\in V$.
\end{lemma}

\begin{proof}
Fix $y$ and let $\xi_i=\mathbf1\{X_{i+1}=y\}-Q_i(X_i,y)$ and $M_t=\sum_{i<t}\xi_i$. The $\xi_i$ are bounded by $1$
and have conditional mean $0$ given $X_0,\dots,X_i$, so they are orthogonal and $\mathbb E M_t^2\le t$. By
Chebyshev's inequality $\Pr[|M_{j^2}|>\zeta j^2]\le\zeta^{-2}j^{-2}$ for every $\zeta>0$, so by the first
Borel--Cantelli lemma $M_{j^2}/j^2\to0$ almost surely. For $j^2\le t<(j+1)^2$ we have $|M_t-M_{j^2}|\le2j+1$, hence
$M_t/t\to0$. Finally $t\mu_t(y)$ and $\sum_{i<t}\mathbf1\{X_{i+1}=y\}$ differ by at most $1$.
\end{proof}

\begin{lemma}\label{lem:limit}
For every strong $\varepsilon$-adversary and every starting distribution, almost surely every limit point $\mu$ of
$(\mu_t)$ satisfies $\mu\bar Q=\mu$ for an admissible matrix $\bar Q$. Consequently, if $\nu(S)\le b$ for every
stationary distribution $\nu$ of every admissible matrix, then almost surely $\limsup_tN_t(S)/t\le b$, and
$\pi^+(S)\le b$.
\end{lemma}

\begin{proof}
Fix an outcome where the conclusion of \Cref{lem:emp} holds, and a sequence $t_j\to\infty$ with $\mu_{t_j}\to\mu$.
Passing to a subsequence we may assume $F_{t_j}\to F$. Then $\mu$ is a probability measure and
$\sum_xF(x,y)=\mu(y)$ for every $y$. Write $F_t(x,\cdot)=(1-\varepsilon)\mu_t(x)P(x,\cdot)+\varepsilon\Gamma_t(x,\cdot)$
with $\Gamma_t(x,\cdot)=\frac1t\sum_{i<t}\mathbf1\{X_i=x\}B_i(x,\cdot)$, a measure of total mass $\mu_t(x)$ supported
in $\supp P(x,\cdot)$. Along $t_j$ it converges to a measure $\Gamma(x,\cdot)$ of total mass $\mu(x)$ supported in
$\supp P(x,\cdot)$, and $F(x,\cdot)=(1-\varepsilon)\mu(x)P(x,\cdot)+\varepsilon\Gamma(x,\cdot)$. Put
$\bar B(x,\cdot)=\Gamma(x,\cdot)/\mu(x)$ if $\mu(x)>0$ and $\bar B(x,\cdot)=P(x,\cdot)$ otherwise. Then
$\bar Q=(1-\varepsilon)P+\varepsilon\bar B$ is admissible and $\mu(x)\bar Q(x,\cdot)=F(x,\cdot)$ for every $x$, so
$\mu\bar Q=\mu$.

For the second statement choose $t_j\to\infty$ with $\mu_{t_j}(S)\to\limsup_t\mu_t(S)$ and, by compactness, with
$\mu_{t_j}$ convergent. Its limit is a stationary distribution of an admissible matrix, so
$\limsup_t\mu_t(S)\le b$. Finally $N_t(S)/t=\mu_t(S)$.
\end{proof}

\section{The complete graph}\label{sec:fail}

\begin{lemma}\label{lem:kn}
For the simple random walk on $K_n$ with $n\ge5$, $\tau=1$ and $\tup(\alpha)=1$ for every $\alpha\in[1/(n-1),1)$.
\end{lemma}

\begin{proof}
From $u$ the law of $X_1$ is uniform on $V\setminus\{u\}$ and $\pi$ is uniform on $V$. The total variation distance
is $\frac12\bigl(\frac1n+(n-1)(\frac1{n-1}-\frac1n)\bigr)=\frac1n\le\frac14$. The ratio $p_1(u,w)/\pi(w)$ is $n/(n-1)$
at every vertex $w\ne u$ and $0$ at $u$, so $s_{\mathrm{up}}(1)=1/(n-1)\le\alpha$. At time $0$ the distance is
$1-1/n>1/4$ and $s_{\mathrm{up}}(0)=n-1>\alpha$, so both times equal $1$.
\end{proof}

\begin{proof}[Proof of \Cref{thm:fail}]
Let the adversary use the point mass $B(x,\cdot)=\delta_v$ at $v$ for $x\ne v$ and $B(v,\cdot)=P(v,\cdot)$. Since
$v$ is a neighbour of every $x\ne v$, $\supp B(x,\cdot)\subseteq\supp P(x,\cdot)$, so
$Q=(1-\varepsilon)P+\varepsilon B$ is admissible. It is irreducible, moves from every $x\ne v$ to $v$ with
probability $a$, and leaves $v$ with probability $1$. If $\nu$ is a stationary distribution of $Q$, then, since
$Q(v,v)=0$ and $Q(x,v)=a$ for $x\ne v$, $\nu(v)=a\sum_{x\ne v}\nu(x)=a(1-\nu(v))$, that is, $\nu(v)=a/(1+a)$. Here
$Q_i=Q$ for all $i$, so $F_t(x,\cdot)=\mu_t(x)Q(x,\cdot)$, and by \Cref{lem:emp} every limit point of $(\mu_t)$ is a
stationary distribution of $Q$. Hence $N_t(\{v\})/t=\mu_t(v)\to a/(1+a)$ almost surely. Since $\varepsilon\le a\le1$,
$a/(1+a)\ge\varepsilon/2$.

Every admissible matrix $Q'$ has $Q'(v,v)=0$ and $Q'(x,v)\le\varepsilon+(1-\varepsilon)/(n-1)=a$ for $x\ne v$. So a
stationary distribution $\nu$ of $Q'$ has $\nu(v)\le a(1-\nu(v))$, that is $\nu(v)\le a/(1+a)$, and
\Cref{lem:limit} gives $\pi^+(\{v\})\le a/(1+a)$.

Let $\alpha\in[1/(n-1),1-1/n]$. Then $\pi(\{v\})=1/n\le1-\alpha$, so \Cref{thm:zu} applies to $S=\{v\}$, and by
\Cref{lem:kn} its bound, with $d=n-1$, is
\[
(1+\alpha)\exp\Bigl(2\varepsilon\sqrt{2(n-1)\ln n}\Bigr)\frac1n=(1+\alpha)\frac{e^{2c\sqrt{2\ln n}}}{n}<\frac{2e^{2c\sqrt{2\ln n}}}{n}.
\]
It is smaller than $\varepsilon/2=c/(2\sqrt{n-1})$ when $c\,n\,e^{-2c\sqrt{2\ln n}}>4\sqrt{n-1}$; the left side grows
like $n^{1-o(1)}$ and the right like $n^{1/2}$.

For $n=65537=2^{16}+1$ and $\varepsilon=2^{-11}$ (so $c=2^{-3}$) we get $a=(33\cdot2^{11}-1)/2^{27}=67583/2^{27}$ and
$a/(1+a)=67583/134285311$, which is $32.98\dots$ times $\pi(\{v\})=1/n$. The exponent in the bound is
$2^{-10}\sqrt{2^{17}\ln n}=2^{-3/2}\sqrt{\ln n}=1.1774\dots$, so the bound is at most $2e^{1.1775}/n<6.50/n$ for
every $\alpha$ in this range and equals $1.25\,e^{1.1774\dots}/n<4.06/n$ for $\alpha=1/4$.
\end{proof}

\begin{proof}[Proof of \Cref{cor:exp}]
For $\kappa\le1/2$ take $c_\kappa=1/\sqrt{20}$ and $C_\kappa=2$, and use \Cref{thm:half} with
$\pi(S)^{1/2}\le\pi(S)^\kappa$. For $\kappa>1/2$, given $c_\kappa$ and $C_\kappa$, the simple random walk on $K_n$ has
$D=n-1$ and $\tau=1$ for $n\ge5$ (\Cref{lem:kn}). With $\varepsilon=\min(c_\kappa,1)/\sqrt{n-1}$, \Cref{thm:fail}
gives $\pi^+(\{v\})\ge\min(c_\kappa,1)/(2\sqrt{n-1})$, which exceeds $C_\kappa n^{-\kappa}$ for large $n$.
\end{proof}

The ``only if'' half of \Cref{cor:exp} also holds with the lazy walk $P_L=(I+P)/2$ on $K_n$ in place of the simple
random walk. Here $D=2(n-1)$. The eigenvalues of $P_L$ other than $1$ equal $r=(n-2)/(2(n-1))$, so
$P_L^t(u,\cdot)-\pi=r^t(\delta_u-\pi)$ and the distance to $\pi$ at time $t$ is $(1-1/n)r^t$. This is
$(n-2)/(2n)>1/4$ at $t=1$ and $(n-2)^2/(4n(n-1))<1/4$ at $t=2$, so $\tau=2$ for $n\ge5$. The lazy walk has a loop at
$v$, so the adversary that always moves to $v$ is allowed at every vertex. Every row of the admissible matrix it uses
gives $v$ mass at least $\varepsilon$, so every stationary distribution of that matrix does too, and as in the proof
of \Cref{thm:fail} the walk spends a fraction at least $\varepsilon$ of its time at $v$. With
$\varepsilon=\min(c_\kappa,1)/\sqrt{4(n-1)}$ this exceeds $C_\kappa n^{-\kappa}$ for large $n$ when $\kappa>1/2$.

\begin{remark}\label{rem:where}
The proof of \Cref{thm:zu} in~\cite{Zu} applies~\cite[Corollary~3.5]{Zu}, which bounds the probability of being in
$S$ at time $t$ under the adversary by $\exp\bigl(2\varepsilon\sqrt{2td\ln(1/p)}\bigr)p$, with $p$ as in
\Cref{sec:intro}. In the example of \Cref{thm:fail} with $n=65537$ and $\varepsilon=2^{-11}$, take $t=1$, a start
$x\ne v$ and $S=\{v\}$, so that $p=1/(n-1)=1/d$. The choice of $k$ in the proof of~\cite[Corollary~3.5]{Zu} then has
$kd\varepsilon=\sqrt{d\ln d/2}$, which is $603$ to the nearest integer. Corollary~3.5 itself fails there. Under the
adversary the walk is at $v$ at time $1$ with probability $a=33/d$ up to a term of order $d^{-1}\varepsilon$, while the
corollary gives $e^{2\varepsilon\sqrt{2d\ln d}}/d<3.3/d$.
\end{remark}

\section{An exponent one half bound}\label{sec:half}

Throughout this section $P$, $\pi$, $D$ and $\tau$ are as in \Cref{thm:half}, and $L^2(\pi)$ is the space of functions
on $V$ with $\|h\|^2=\sum_x\pi(x)h(x)^2$. For a matrix $K$ on $V$ write $Kh(x)=\sum_yK(x,y)h(y)$, and let $K^*$ be its
adjoint in $L^2(\pi)$, $K^*g(y)=\sum_x\pi(x)g(x)K(x,y)/\pi(y)$. If $\nu$ is a measure with density $f=\nu/\pi$, then
$\nu K$ has density $K^*f$. Reversibility says that $P^*=P$.

\begin{lemma}\label{lem:lr}
If $Q$ is admissible and $\varepsilon^2D\tau\le1/20$, then $\|(Q^\tau -P^\tau)h\|\le\frac14\|h\|$ for every $h$.
\end{lemma}

\begin{proof}
On $\supp P$ write $Q(x,y)/P(x,y)=1+\varepsilon(U(x,y)-1)$ with $U=B/P$. Then $0\le U\le D$ and
$\sum_yP(x,y)U(x,y)=1$, so
\[
\sum_yP(x,y)\bigl(1+\varepsilon(U(x,y)-1)\bigr)^2=1+\varepsilon^2\Bigl(\sum_yP(x,y)U(x,y)^2-1\Bigr)\le1+\varepsilon^2(D-1).
\]
Let $(Y_i)$ be the chain $P$ started at $x$, and $L=\prod_{i<\tau}Q(Y_i,Y_{i+1})/P(Y_i,Y_{i+1})$. Since $Q(x,y)>0$ only
where $P(x,y)>0$, $Q^\tau h(x)=\mathbb E_x[L\,h(Y_\tau)]$ and $\mathbb E_x L=1$. Conditioning on $Y_0,\dots,Y_{\tau -1}$
and using the display $\tau$ times gives $\mathbb E_xL^2\le(1+\varepsilon^2D)^\tau\le e^{\varepsilon^2D\tau}$, so
$\mathbb E_x(L-1)^2\le e^{1/20}-1<1/16$. By Cauchy--Schwarz,
\[
\bigl|(Q^\tau -P^\tau)h(x)\bigr|=\bigl|\mathbb E_x[(L-1)h(Y_\tau)]\bigr|\le\tfrac14\bigl(P^\tau (h^2)(x)\bigr)^{1/2}.
\]
Squaring, multiplying by $\pi(x)$ and summing, with $\pi P^\tau=\pi$, gives the claim.
\end{proof}

\begin{lemma}\label{lem:gap}
Let $A=P^\tau -\Pi$, where $\Pi(x,y)=\pi(y)$. Then $\|Ah\|\le\frac12\|h\|$ for every $h$.
\end{lemma}

\begin{proof}
By the definition of $\tau$, $\sum_y|A(x,y)|\le\frac12$ for every $x$. By reversibility
$\pi(x)P^\tau (x,y)=\pi(y)P^\tau (y,x)$, so
\[
\sum_x\pi(x)|A(x,y)|=\pi(y)\sum_x|P^\tau (y,x)-\pi(x)|\le\tfrac12\pi(y).
\]
By Cauchy--Schwarz, $|Ah(x)|^2\le\frac12\sum_y|A(x,y)|h(y)^2$, and summing against $\pi$ gives
$\|Ah\|^2\le\frac12\sum_yh(y)^2\sum_x\pi(x)|A(x,y)|\le\frac14\|h\|^2$.
\end{proof}

\begin{proposition}\label{prop:stat}
Let $Q$ be admissible with $\varepsilon^2D\tau\le1/20$, and let $\nu$ be a stationary distribution of $Q$. Then
$\nu(S)\le2\pi(S)^{1/2}$ for every $S\subseteq V$.
\end{proposition}

\begin{proof}
Let $f=\nu/\pi$. Since $\nu Q^\tau=\nu$, $f=(Q^\tau)^*f$. Since $\sum_x\pi(x)f(x)=1$, $\Pi^*f$ is the constant
function $\mathbf1$; also $(P^\tau)^*\mathbf1=\mathbf1$ and $\Pi^*\mathbf1=\mathbf1$, so $A^*\mathbf1=0$ and
\[
f-\mathbf1=(Q^\tau -P^\tau)^*f+(P^\tau)^*f-\Pi^*f=(Q^\tau -P^\tau)^*f+A^*(f-\mathbf1).
\]
Adjoints have the same norm, so by \Cref{lem:lr,lem:gap}, $\|f-\mathbf1\|\le\frac14\|f\|+\frac12\|f-\mathbf1\|$, that is
$\|f-\mathbf1\|\le\frac12\|f\|$, and then $\|f\|\le1+\frac12\|f\|$ gives $\|f\|\le2$. By Cauchy--Schwarz,
$\nu(S)=\sum_{x\in S}\pi(x)f(x)\le\|f\|\pi(S)^{1/2}\le2\pi(S)^{1/2}$.
\end{proof}

\begin{proof}[Proof of \Cref{thm:half}]
By \Cref{prop:stat}, every stationary distribution $\nu$ of every admissible matrix has $\nu(S)\le2\pi(S)^{1/2}$, and
\Cref{lem:limit} gives the claim.
\end{proof}

The next bound needs neither reversibility nor a lower bound on the entries of $P$.

\begin{proposition}\label{prop:tv}
Let $P$ be an irreducible transition matrix with stationary distribution $\pi$ and mixing time $\tau<\infty$. If
$Q=(1-\varepsilon)P+\varepsilon B$ for a transition matrix $B$ and $\nu$ is a stationary distribution of $Q$, then
$\|\nu-\pi\|_{\mathrm{TV}}\le2\varepsilon\tau$. Consequently $\pi^+(S)\le\pi(S)+2\varepsilon\tau$ for every
$S\subseteq V$.
\end{proposition}

\begin{proof}
For probability measures $\rho$ on $V$, $\|\rho Q-\rho P\|_{\mathrm{TV}}=\varepsilon\|\rho B-\rho
P\|_{\mathrm{TV}}\le\varepsilon$. Writing $\rho Q^\tau-\rho P^\tau=\sum_{j<\tau}(\rho Q^{j}Q-\rho
Q^{j}P)P^{\tau-1-j}$ and using that $P$ does not increase total variation gives
$\|\rho Q^\tau-\rho P^\tau\|_{\mathrm{TV}}\le\varepsilon\tau$. If $\eta$ is a signed measure of total mass $0$ and
$s=\|\eta\|_{\mathrm{TV}}>0$, then $\eta^+/s$ and $\eta^-/s$ are probability measures whose difference is $\eta/s$; by
convexity $\|\eta P^\tau\|_{\mathrm{TV}}\le s\max_{x,y}\|P^\tau(x,\cdot)-P^\tau(y,\cdot)\|_{\mathrm{TV}}
\le\frac12\|\eta\|_{\mathrm{TV}}$, by the triangle inequality and the definition of $\tau$. Applying both to
$\nu-\pi=(\nu Q^\tau-\nu P^\tau)+(\nu-\pi)P^\tau$ gives
$\|\nu-\pi\|_{\mathrm{TV}}\le\varepsilon\tau+\frac12\|\nu-\pi\|_{\mathrm{TV}}$. The second statement follows from the
first and \Cref{lem:limit}.
\end{proof}

\begin{remark}\label{rem:mix}
The first step of this proof is the estimate behind~\cite[Proposition~2.12]{Zu}.
\end{remark}

\section{Comparison with other bounds}\label{sec:compare}

For a graph with minimum degree at least $\gamma n$ and $\varepsilon\le\beta/n$, Haslegrave, Sauerwald and
Sylvester~\cite[Proposition~4.7]{HSS} prove that a stationary distribution of one fixed admissible matrix gives each
vertex $u$ probability at most $(1+\beta)\gamma^{-2}\pi(u)$. By \Cref{lem:limit} a bound for every admissible matrix
is a bound on $\pi^+$. On $K_n$, with $\gamma=1-1/n$ and $\beta=\varepsilon n$, this gives
$\pi^+(\{v\})\le(1+\varepsilon n)n/(n-1)^2$, which is of order $\varepsilon$ for the $\varepsilon$ of
\Cref{thm:fail}. There \Cref{thm:fail} gives the exact value $a/(1+a)$.

\Cref{prop:tv} gives $\pi^+(S)\le\pi(S)+2\varepsilon\tau$ for every chain. If $2\varepsilon\tau<1$, then
$2\pi(S)^{1/2}<\pi(S)+2\varepsilon\tau$ exactly when $\pi(S)<\bigl(1-\sqrt{1-2\varepsilon\tau}\bigr)^2$, which is
$(\varepsilon\tau)^2(1+O(\varepsilon\tau))$; under the hypothesis of \Cref{thm:half}, $(\varepsilon\tau)^2$ is at most
$\tau/(20D)$. If $2\varepsilon\tau\ge1$, the bound of \Cref{prop:tv} is at least $1$, and \Cref{thm:half} is the
smaller one for every $S$ with $\pi(S)<1/4$. On $K_n$ the two bounds are of the same order for a single vertex.

Let $d(u)$ be the degree of a vertex $u$. Haslegrave, Sauerwald and Sylvester~\cite[Section~4.3]{HSS} call a matrix $Q$
supported on the edges of a graph a $(\theta,\Theta)$-simple walk if $\theta/d(u)\le Q(u,w)\le\Theta/d(u)$ for every
edge $uw$. They prove that its stationary distribution gives each vertex $x$ probability at most
$2\Theta^{\tau_2}\pi(x)$ \cite[Proposition~4.8]{HSS}. Here $\tau_2=\min\{t_\infty,3\ln n/|\ln\lambda_*|\}$. The time
$t_\infty$, defined in \cite[Section~2, equation~(2)]{HSS}, is the least $t$ with $|\tilde p_t(x,y)/\pi(y)-1|<1/e$ for all $x,y$, where $\tilde p_t$ is the law at time
$t$ of the lazy walk $(I+P)/2$, and $\lambda_*$ is the largest absolute value of an eigenvalue other than $1$ of the
simple random walk. On a graph with degrees at most $d$ an admissible matrix satisfies
$(1-\varepsilon)/d(u)\le Q(u,w)\le\varepsilon+(1-\varepsilon)/d(u)\le(1+\varepsilon d)/d(u)$, as in the proof
of~\cite[Proposition~4.7]{HSS}. So it is a $(1-\varepsilon,1+\varepsilon d)$-simple walk, and by \Cref{lem:limit}
$\pi^+(\{x\})\le2(1+\varepsilon d)^{\tau_2}\pi(x)$. This bound is the stronger one when $\varepsilon d\tau_2$ is
bounded. Where both apply, \Cref{thm:half} gives the smaller bound at $x$ exactly when
$(1+\varepsilon d)^{\tau_2}>\pi(x)^{-1/2}$. On $K_n$ with $\varepsilon=c/\sqrt{n-1}$ and $c\le1/\sqrt{20}$ we have
$t_\infty\ge2$, since the lazy walk stays put at time $1$ with probability $1/2$, and $3\ln n/\ln(n-1)>3$. So
$\tau_2\ge2$, the bound of~\cite{HSS} at a vertex is at least $2(1+c\sqrt{n-1})^2/n>2c^2(n-1)/n$, and \Cref{thm:half}
gives $2n^{-1/2}$.

The rook's graph $K_q\square K_q$ has the pairs $(i,j)$ with $1\le i,j\le q$ as vertices, and joins two pairs when
they agree in exactly one coordinate. It has $n=q^2$ vertices, all of degree $d=2(q-1)$, it is connected, and for
$q\ge3$ it contains triangles.

\begin{lemma}\label{lem:rook}
For $q\ge3$ the simple random walk on $K_q\square K_q$ has $\tau\le4$.
\end{lemma}

\begin{proof}
A step of the walk picks one of the two coordinates with probability $1/2$ and moves it to a uniform one of the other
$q-1$ values. Given that $j$ of the first $t$ steps move the first coordinate, the two coordinates are independent, and
they are the simple random walk on $K_q$ after $j$ and after $t-j$ steps. After $m\ge1$ steps that walk is at its
starting value with probability $1/q+(1-1/q)(-1/(q-1))^m$ and is equally likely to be at each other value, so its
total variation distance to the uniform law is $(1-1/q)(q-1)^{-m}\le1/q$. The distance of a product of two laws from
the uniform law on pairs is at most the sum of the two distances, and with probability $2^{1-t}$ one coordinate is
never moved. So the distance to $\pi$ at time $t$ is at most $2^{1-t}+2/q$, which is at most $1/4$ at $t=4$ when
$q\ge16$. For $3\le q\le15$ the exact distance, computed by the script of the data section, is at most $1/4$ at
$t=3$.
\end{proof}

The same computation gives $\tau=2$ for $3\le q\le7$ and $\tau=3$ for $8\le q\le1000$.

\begin{proposition}\label{prop:rook}
Let $q\ge41$, let $\tau$ be the mixing time of the simple random walk on $K_q\square K_q$, and let
$\varepsilon=1/\sqrt{20d\tau}$ with $d=2(q-1)$. Then \Cref{thm:half} applies with $D=d$ and gives
$\pi^+(\{x\})\le2/q$ for every vertex $x$. At every vertex $x$:
\begin{enumerate}[label=(\roman*)]
\item the bound $\pi(\{x\})+2\varepsilon\tau$ of \Cref{prop:tv} is larger than $2/q$;
\item the bound $(1+\beta)\gamma^{-2}\pi(x)$ of~\cite[Proposition~4.7]{HSS}, with $\gamma=d/n$, is larger than $1/4$;
\item the bound $2(1+\varepsilon d)^{\tau_2}\pi(x)$ from~\cite[Proposition~4.8]{HSS} is larger than $1$;
\item $2d\varepsilon>1$, so \cite[Theorem~3.3]{Zu} with $k=2$ does not apply.
\end{enumerate}
\end{proposition}

\begin{proof}
The walk is irreducible and aperiodic, reversible with respect to the uniform law, and its positive entries equal
$1/d$. By the choice of $\varepsilon$, $\varepsilon^2d\tau=1/20$, and \Cref{thm:half} gives $2\pi(\{x\})^{1/2}=2/q$.
By \Cref{lem:rook}, $1\le\tau\le4$.

(i) $2\varepsilon\tau=\sqrt{\tau/(5d)}=\sqrt{\tau/(10(q-1))}\ge1/\sqrt{10(q-1)}$, and this exceeds $2/q$ because
$q^2>40(q-1)$ for $q\ge39$.

(ii) $\gamma^{-2}\pi(x)=q^4/(4(q-1)^2)\cdot q^{-2}=q^2/(4(q-1)^2)>1/4$, and $1+\beta\ge1$.

(iii) Since $q-1\ge40\ge10\tau$, $\varepsilon d=\sqrt{d/(20\tau)}=\sqrt{(q-1)/(10\tau)}\ge1$. The lazy walk stays at $x$
for its first $t$ steps with probability $2^{-t}$, so $\tilde p_t(x,x)/\pi(x)-1\ge n2^{-t}-1$, and therefore
$2^{t_\infty}>n/(1+e^{-1})$. The eigenvalues of the simple random walk other than $1$ are
$(q-2)/(2(q-1))$ and $-1/(q-1)$, so $\lambda_*=(q-2)/(2(q-1))\ge1/3$ and $3\ln n/|\ln\lambda_*|\ge3\ln n/\ln3>\ln n/\ln2$.
Hence $2^{\tau_2}>n/(1+e^{-1})$, and $2(1+\varepsilon d)^{\tau_2}\pi(x)\ge2\cdot2^{\tau_2}/n>2/(1+e^{-1})>1$.

(iv) $2d\varepsilon=2\sqrt{d/(20\tau)}\ge2$ by (iii).
\end{proof}

In the other direction, a controller that fixes one matrix $B$ can raise the stationary probability of a vertex
$x$ to at least $\pi(x)^{1-\varepsilon}/(16\,\tilde t)$, where $\tilde t$ is the total variation mixing time of the
lazy walk~\cite[Theorem~4.3]{HSS}. Azar et al.~\cite{ABKLP} proved $\pi(x)^{1-\Omega(\varepsilon)}$ for regular graphs of bounded
degree and conjectured $\pi(x)^{1-\varepsilon}$ for every graph; the conjecture fails for $K_2$, as recorded
in~\cite[Section~4]{HSS}. These are lower bounds for $\pi^+(\{x\})$.

\section*{Acknowledgments and funding}

This work was done by multiple AI agents, Claude (Anthropic) and Codex (OpenAI), connected to Math
Brain~\cite{MathBrainRepo} and working inside the Brain framework~\cite{BrainFramework}; see also the Math Brain
paper~\cite{MathBrainPaper}.

The research presented in this paper was supported by the European Research Council (ERC) under the European
Union's Horizon Europe research and innovation programme (grant agreement No.\ 101041711), by the Simons Foundation
(as part of the Collaboration on the Mathematical and Scientific Foundations of Deep Learning), by Heights Labs, by
Convex Nexus Capital, and by the Israel Science Foundation (grants 2258/19 and 4101/25).

\section*{Declaration of competing interest}

The author holds equity in Heights Labs and Convex Nexus Capital and has a management role at Convex Nexus Capital.

\section*{Declaration of generative AI and AI-assisted technologies}

The results of \Cref{sec:emp,sec:fail,sec:half,sec:compare} were developed by OpenAI Codex and Anthropic Claude
agents working under the author's direction with the shared memory of Math Brain \cite{MathBrainRepo,MathBrainPaper}
and the Brain framework \cite{BrainFramework}. The agents also drafted the text and wrote the exact-arithmetic script
of the data section. The model, the definitions and \Cref{thm:zu} are Zuckerman's \cite{Zu}, and the first step of
\Cref{prop:tv} is his Proposition~2.12. The controller model is due to Azar, Broder, Karlin, Linial and Phillips
\cite{ABKLP}. The bounds compared in \Cref{sec:compare} are Propositions~4.7 and~4.8 of Haslegrave, Sauerwald and
Sylvester \cite{HSS}. The author directed the work and is responsible for the content of the paper and for decisions
concerning its dissemination.

\section*{Data and code availability}

A Python script that recomputes the numbers in the proof of \Cref{thm:fail} with exact rational arithmetic, checks the
stationary mass against the full matrix for small $n$, and computes the mixing time of the rook's graph used in
\Cref{lem:rook}, is available at \url{https://github.com/yspennstate/strong-adversary-large-degree}.

\appendix
\section{Notation}

\begin{center}
\small
\begin{tabular}{lp{0.5\textwidth}l}
Symbol & Meaning & First use\\[2pt]
$P$, $V$, $\pi$ & transition matrix, state space, stationary distribution & \Cref{sec:intro}\\
$\varepsilon$, $B_i$, $X_i$ & adversary's share, its matrix and the walk at time $i$ & \Cref{sec:intro}\\
$b$ & a bound on the stationary probabilities of $S$ & \Cref{lem:limit}\\
$Q$, $Q_i$ & an admissible matrix; the one used at time $i$ & \Cref{sec:intro}\\
$S$, $N_t(S)$, $\pi^+(S)$ & a set of states; visits to $S$ before time $t$; upper stationary probability & \Cref{sec:intro}\\
$\tau$ & mixing time at total variation $1/4$ & \Cref{sec:intro}\\
$p_t(u,\cdot)$ & law at time $t$ of the walk without the adversary started at $u$ & \Cref{sec:intro}\\
$s_{\mathrm{up}}$, $\alpha$, $\tup(\alpha)$ & upper separation, its threshold, upper separation time & \Cref{sec:intro}\\
$d$, $G$ & bound on the degrees; a graph & \Cref{thm:zu}\\
$K_n$, $n$, $v$ & complete graph on $n$ vertices; a fixed vertex & \Cref{sec:intro}\\
$c$ & constant in the size of $\varepsilon$ on $K_n$ & \Cref{sec:intro}\\
$p$, $k$ & probability of $S$ at time $t$ without the adversary; exponent in~\cite[Theorem~3.3]{Zu} & \Cref{sec:intro}\\
$a$ & probability that the walk under the adversary of \Cref{thm:fail} enters $v$ & \Cref{thm:fail}\\
$D$ & reciprocal of the smallest positive entry of $P$ & \Cref{sec:intro}\\
$Y_i$, $\mathbb E_x$, $L$ & chain $P$ started at $x$, its expectation, likelihood ratio & \Cref{sec:intro}\\
$h$, $\|h\|$ & a function on $V$; its norm in $L^2(\pi)$ & \Cref{sec:intro}\\
$\Pi$ & the matrix with rows $\pi$ & \Cref{sec:intro}\\
$\nu$, $f$, $\mathbf 1$ & a stationary distribution of an admissible matrix; its density; the constant function & \Cref{sec:intro}\\
$\kappa$, $c_\kappa$, $C_\kappa$ & exponent and constants of the corollary & \Cref{cor:exp}\\
$q$, $K_q\square K_q$ & side and rook's graph & \Cref{sec:intro}\\
$\mu_t$, $F_t$ & empirical occupation and flow measures & \Cref{sec:emp}\\
$\xi_i$, $M_t$, $\zeta$ & martingale differences, their sum, a positive constant & \Cref{lem:emp}\\
$t_j$, $\mu$, $F$ & a sequence of times; limits of $\mu_{t_j}$ and $F_{t_j}$ & \Cref{lem:limit}\\
$\Gamma_t$, $\Gamma$, $\bar B$, $\bar Q$ & adversary measure, its limit, and the matrices built from it & \Cref{lem:limit}\\
$\delta_v$ & point mass at $v$ & \Cref{sec:fail}\\
$P_L$, $r$ & lazy walk on $K_n$; its eigenvalue other than $1$ & \Cref{sec:fail}\\
$K^*$ & adjoint of $K$ in $L^2(\pi)$ & \Cref{sec:half}\\
$U$ & the ratio $B/P$ & \Cref{lem:lr}\\
$A$ & $P^\tau -\Pi$ & \Cref{lem:gap}\\
$\rho$, $\eta$, $s$ & probability measure; signed measure of mass $0$; its total variation & \Cref{prop:tv}\\
$\beta$, $\gamma$ & constants in the dense-graph bound of~\cite{HSS} & \Cref{sec:compare}\\
$d(u)$, $u$, $w$ & degree of the vertex $u$; vertices & \Cref{sec:compare}\\
$\theta$, $\Theta$ & constants of a simple walk in~\cite{HSS} & \Cref{sec:compare}\\
$\tau_2$, $t_\infty$, $\tilde p_t$, $\lambda_*$ & quantities in~\cite[Proposition~4.8]{HSS} & \Cref{sec:compare}\\
$m$, $j$ & numbers of steps in the proof of \Cref{lem:rook} & \Cref{lem:rook}\\
$\tilde t$ & total variation mixing time of the lazy walk & \Cref{sec:compare}\\
\end{tabular}
\end{center}

\begin{thebibliography}{9}

\bibitem{ABKLP}
Y.~Azar, A.~Z.~Broder, A.~R.~Karlin, N.~Linial and S.~J.~Phillips, \emph{Biased random walks}, Combinatorica
\textbf{16} (1996), 1--18. \doi{10.1007/BF01300124}

\bibitem{HSS}
J.~Haslegrave, T.~Sauerwald and J.~Sylvester, \emph{Time dependent biased random walks}, ACM Trans.\ Algorithms
\textbf{18} (2022), no.~2, 1--30. \doi{10.1145/3498848}. Statement numbers refer to arXiv:2006.02475v2.

\bibitem{MathBrainRepo}
Y.~Shmalo, \emph{Math Brain: public snapshot, code and tests}, 2026.
\url{https://github.com/yspennstate/math-brain-public}

\bibitem{MathBrainPaper}
Y.~Shmalo, \emph{Math Brain: shared, auditable memory for multi-agent mathematical research}, October 2026.
\url{https://github.com/yspennstate/math-brain-public/tree/main/papers/math-brain-2026-10-24}

\bibitem{BrainFramework}
Y.~Shmalo, \emph{Brain Framework}, public software repository, 2026.
\url{https://github.com/yspennstate/brain-framework-public}

\bibitem{Zu}
D.~Zuckerman, \emph{Markov chain robustness}, in: 17th Innovations in Theoretical Computer Science Conference
(ITCS 2026), LIPIcs \textbf{362}, Paper~118. \doi{10.4230/LIPIcs.ITCS.2026.118}

\end{thebibliography}

\end{document}
